\section{Dimensions of labels on the bad copies}
\label{sec:bad_copy_label_dimension}

The auxiliary estimate in Section~7 uses a bound on the dimensions of labels
occurring under permutations of the bad copies. The good copies may remain
in the representation space, but their multiplicity does not enter this bound.

\begin{theorem}[Dimension of an occurring bad-copy label]
\label{thm:bad_copy_label_dimension}
\lean{TensorPower.groupHom₂_labelProj_dim_le}
\leanok
Let $C$ be a finite coordinate set with $d=|C|$, and split the $k$ copy
coordinates into specified groups of sizes $m$ and $r$, where $m+r=k$.
Let $S_r$ act on $\mathbb C^{C^k}$ by permuting the second group and fixing
the first group. If its central projection $P_\beta$ for an irreducible
label $\beta$ is nonzero, then
\[
 d_\beta\le d^r.
\]
The statement includes $m=0$, $r=0$, and $C=\varnothing$.
\end{theorem}
\begin{proof}
\leanok
Write $x=(x_g,x_b)$ and $y=(y_g,y_b)$ in the specified split of the
coordinates. The central projection has entries
\[
 P_\beta(x,y)=\mathbf 1_{x_g=y_g}\,P^{(r)}_\beta(x_b,y_b),
\]
where $P^{(r)}_\beta$ is the same label projection for the actual
permutation action on $\mathbb C^{C^r}$. Thus nonvanishing of $P_\beta$
implies nonvanishing of $P^{(r)}_\beta$. Its range is invariant under
$S_r$, is nonzero, and is contained in the $\beta$-isotypic component.
Consequently its dimension is at least $d_\beta$ and at most the ambient
dimension $d^r$.
\end{proof}

This proves the bad-label dimension bound used in
\cite[Section~7, equation \textup{(comparator:good-auxiliary)}]{OpenAI2026AreaLaw}.
The preservation of whole-copy auxiliary labels by the physical excitation
projection, and the resulting good-auxiliary estimate, are separate assertions.
